\chapter{Caffeine}
\index{Caffeine}

\section*{Definition and Overview}
Caffeine is a naturally occurring stimulant found in coffee, tea, cocoa, and many soft drinks and energy drinks. It belongs to a class of substances called methylxanthines and works by stimulating the central nervous system, making people feel more alert and less tired. Caffeine is the most widely consumed psychoactive substance in the world.

At normal dietary levels, caffeine is generally safe for healthy adults. However, it can cause side effects in sensitive people or at high doses, it interacts with sleep and anxiety, and some people become dependent on it, experiencing withdrawal symptoms when they stop. Understanding how much caffeine you consume and how it affects you is central to using it sensibly.

\section*{Epidemiology}
Caffeine is consumed by the majority of adults worldwide, most commonly as coffee and tea. In many countries, roughly three in four adults drink coffee regularly, and average intake is a few hundred milligrams a day, equivalent to two to four cups. Teenagers and young adults obtain more of their caffeine from soft drinks and energy drinks, and energy-drink consumption has grown sharply over the past two decades. Pregnant women are advised to limit caffeine, although many continue to consume modest amounts. Variability in intake reflects differences in beverage choices, portion sizes, and cultural habits.

\section*{Aetiology and Pathophysiology}
Caffeine acts mainly by blocking the effects of adenosine, a chemical that builds up in the brain during the day and promotes sleepiness:
\begin{itemize}
    \item \textbf{Adenosine blockade:} by blocking adenosine receptors, caffeine reduces feelings of tiredness and maintains alertness
    \item \textbf{Stimulation:} it also raises levels of other brain chemicals and increases heart rate, blood pressure, and the release of stress hormones
    \item \textbf{Metabolism:} caffeine is broken down by the liver, and its effects last for several hours; breakdown is slower in pregnancy, in newborns, and in people taking certain medicines
    \item \textbf{Tolerance and withdrawal:} regular use produces tolerance, and sudden cessation can cause headache, tiredness, and irritability
    \item \textbf{Sources and dose:} a cup of coffee contains roughly 80--100 mg, tea 30--50 mg, and a can of energy drink 80--150 mg of caffeine
\end{itemize}

\section*{Clinical Features}
The effects of caffeine vary with the dose and with individual sensitivity:
\begin{itemize}
    \item \textbf{Alertness:} increased attention, improved concentration, and reduced fatigue at normal doses
    \item \textbf{Overstimulation:} restlessness, nervousness, irritability, and difficulty sleeping at higher doses
    \item \textbf{Palpitations:} a rapid, irregular, or forceful heartbeat in some people
    \item \textbf{Neurological effects:} headache, dizziness, and tremor in sensitive individuals
    \item \textbf{Gastrointestinal effects:} stomach upset, heartburn, and increased urination
    \item \textbf{Anxiety:} worsening of anxiety and panic symptoms, especially in people already prone to them
    \item \textbf{Withdrawal:} headache, tiredness, low mood, and poor concentration after stopping regular use, usually beginning within a day and settling over a few days
\end{itemize}

\section*{Differential Diagnosis}
Symptoms caused by caffeine can mimic other conditions:
\begin{itemize}
    \item Anxiety and panic disorder, which caffeine can trigger or worsen
    \item Heart conditions causing palpitations or an irregular heartbeat
    \item Migraine and other headache disorders
    \item Sleep disorders such as insomnia or fragmented sleep
    \item Medication side effects or withdrawal from other substances
    \item Medical causes of tremor, such as an overactive thyroid
\end{itemize}

\section*{When to Seek Medical Care}
Most caffeine-related symptoms settle when intake is reduced, but some symptoms need urgent assessment.

\textbf{Contact your local emergency services for:}
\begin{itemize}
    \item Chest pain or a very fast, irregular, or forceful heartbeat
    \item Fainting, collapse, or seizures
    \item Severe agitation, confusion, or hallucinations after a large intake of caffeine
    \item Symptoms suggesting an allergic reaction, such as difficulty breathing or swelling of the face
\end{itemize}

Otherwise, people should see a doctor if anxiety, palpitations, headaches, or sleep problems are persistent or troublesome.

\section*{Diagnosis and Investigations}
There is no specific test for caffeine use; the diagnosis is made from the history and by excluding other causes:
\begin{itemize}
    \item \textbf{History:} a careful record of how much coffee, tea, cola, chocolate, and energy drinks is consumed daily
    \item \textbf{Examination:} assessment of heart rate, blood pressure, and signs of overstimulation such as tremor
    \item \textbf{Blood tests:} thyroid function, blood glucose, and electrolytes may be checked if symptoms suggest another cause
    \item \textbf{Electrocardiogram:} a heart tracing may be arranged if palpitations are reported
    \item \textbf{Medication review:} checking for other stimulants or medicines that interact with caffeine
\end{itemize}

\section*{Management}
Management usually involves adjusting intake rather than stopping caffeine altogether:
\begin{itemize}
    \item \textbf{Reduce intake:} cutting back gradually to a moderate level of around 200--400 mg a day, or two to four cups of coffee, suits most people
    \item \textbf{Taper rather than stop suddenly:} reducing gradually over one to two weeks prevents withdrawal headaches
    \item \textbf{Timing:} avoiding caffeine in the afternoon and evening to protect sleep
    \item \textbf{Substitutes:} choosing decaffeinated coffee or tea, or smaller serves
    \item \textbf{Symptom relief:} paracetamol for withdrawal headaches while reducing intake
    \item \textbf{Address the cause:} if caffeine is used to mask fatigue, improving sleep and treating any underlying anxiety is more effective in the long term
    \item \textbf{Special groups:} pregnant women are advised to keep caffeine below about 200 mg a day
\end{itemize}

\section*{Complications}
\begin{itemize}
    \item Dependence, with withdrawal symptoms on stopping
    \item Worsening of anxiety, panic attacks, and insomnia
    \item Palpitations and, very rarely, serious heart rhythm disturbances after very high doses
    \item Sleep disruption, which affects mood, energy, and concentration
    \item Heartburn and stomach irritation in sensitive people
    \item Accidental overdose from concentrated caffeine powders or caffeine pills, which is rare but dangerous
\end{itemize}

\section*{Prognosis and Outcomes}
For most people, caffeine is safe, and its side effects resolve quickly once intake is reduced. Withdrawal symptoms settle within a few days of cutting down. People who are sensitive, anxious, or prone to heart rhythm problems usually feel noticeably better after reducing or retiming their caffeine intake. There is no evidence that moderate caffeine consumption causes lasting harm in otherwise healthy adults.

\section*{Prevention and Self-Care}
\begin{itemize}
    \item Stay within moderate amounts, roughly two to four cups of coffee a day, and avoid very large single servings
    \item Be careful with energy drinks, which can contain large amounts of caffeine alongside other stimulants
    \item Keep track of caffeine from all sources, including tea, cola, chocolate, and some medicines
    \item Avoid caffeine in the late afternoon and evening if you struggle to sleep
    \item If you decide to cut down, reduce gradually to avoid withdrawal headaches
    \item Limit caffeine during pregnancy and discuss it with your doctor or midwife
    \item Seek help if anxiety or poor sleep is persistent, as caffeine is often a mask rather than a cure
\end{itemize}

\section*{Sources and Further Reading}
The following sources provide reliable, accessible information on caffeine and its effects on health:
\begin{itemize}
    \item NHS. \textit{Caffeine and your health.} https://www.nhs.uk/live-well/eat-well/
    \item MedlinePlus. \textit{Caffeine.} https://medlineplus.gov/caffeine.html
    \item Mayo Clinic. \textit{Caffeine: how much is too much?} https://www.mayoclinic.org/
    \item Food Standards Australia New Zealand. \textit{Caffeine and energy drinks.} https://www.foodstandards.gov.au/
    \item MSD Manual. \textit{Caffeine.} https://www.msdmanuals.com/
\end{itemize}
